% Generated draft: requires mathematical and source review.
\documentclass[11pt]{article}
\usepackage[margin=25mm]{geometry}
\usepackage{fontspec}
\setmainfont{Latin Modern Roman}
\usepackage{amsmath,amssymb,mathtools,unicode-math}
\setmathfont{Latin Modern Math}
\usepackage{enumitem,needspace,xurl,hyperref}
\setlength{\parindent}{0pt}
\setlength{\parskip}{5pt}
\setlength{\emergencystretch}{3em}
\newcommand{\N}{\mathbb{N}}
\newcommand{\Z}{\mathbb{Z}}
\newcommand{\Q}{\mathbb{Q}}
\newcommand{\R}{\mathbb{R}}
\newcommand{\C}{\mathbb{C}}
\newcommand{\F}{\mathbb{F}}
\newcommand{\A}{\mathbb{A}}
\newcommand{\PP}{\mathbb{P}}
\newcommand{\E}{\mathbb{E}}
\newcommand{\eps}{\varepsilon}
\newcommand{\dd}{\,\mathrm{d}}
\begin{document}
\section*{UnsolvedMath ID: 701593}
\textbf{Source problem code:} catalogue entry
\subsection{Definitions and mathematical statement}
Write \(\mathbb{N}_0=\{0,1,2,\ldots\}\). Let \(G=\langle t_1,t_2,t_3\rangle\) be a group of permutations of \(\mathbb{Z}\), where the three specified generators are class transpositions in the sense of Kourovka Notebook Problem 17.57. That definition is referenced but not reproduced in the supplied record; consequently, the precise permitted generators and their permutation formulas remain dependent on the missing definition.

For \(i\in\{1,2,3\}\), denote the two residue classes interchanged by \(t_i\) by \(a_i+u_i\mathbb{Z}\) and \(b_i+v_i\mathbb{Z}\), where \(u_i,v_i\) are their positive moduli, and put
\[
 m=\operatorname{lcm}(u_1,v_1,u_2,v_2,u_3,v_3).
\]
For each \(S\subseteq\{0,1,\ldots,m-1\}\), define
\[
 U_S=\bigcup_{r\in S}(r+m\mathbb{Z}).
\]
Assume that
\[
 \forall S\subseteq\{0,1,\ldots,m-1\},\qquad
 \bigl[(\forall g\in G)\;g(U_S)=U_S\bigr]
 \Longrightarrow
 \bigl[S=\varnothing\ \text{or}\ S=\{0,1,\ldots,m-1\}\bigr].
\]
Assume also that the integers \(0,1,\ldots,42\) belong to one orbit of the action on \(\mathbb{Z}\); equivalently,
\[
 (\forall k\in\{0,1,\ldots,42\})(\exists g\in G)\quad g(0)=k.
\]
The question asks whether the action of \(G\) on \(\mathbb{N}_0\) must be transitive, that is, whether
\[
 (\forall a,b\in\mathbb{N}_0)(\exists g\in G)\quad g(a)=b.
\]
Calling this a restricted action requires \(g(\mathbb{N}_0)=\mathbb{N}_0\) for every \(g\in G\). The supplied statement presupposes such an action, but without the referenced definition of class transpositions this invariance cannot be checked from the supplied record. The displayed conclusion expresses the intended transitivity question subject to that unresolved qualification.

The source further asserts that the bound \(42\) cannot be decreased: in its notation, the finite group
\[
 H=\left\langle
 \tau_{0(2),1(2)},\,
 \tau_{0(3),2(3)},\,
 \tau_{0(7),6(7)}
 \right\rangle
\]
acts transitively on \(\{0,1,\ldots,41\}\) and on the set of residue classes modulo \(42\). Here \(a(d)\) denotes the residue class \(a+d\mathbb{Z}\), while the exact permutation convention for \(\tau\) remains the one in the missing referenced definition. These are supplied source assertions, retained without independent verification.
\subsection{Short English statement}
S. Kohl asks whether three class transpositions must act transitively on the nonnegative integers when they preserve no nontrivial union of residue classes modulo the specified least common multiple and place \(0,\ldots,42\) in one orbit. The source supplies a finite example to assert that the bound cannot be lowered. Precision still requires the referenced definition of class transpositions and confirmation that their action preserves \(\mathbb{N}_0\).
\subsection{Sources}
\begin{itemize}
\item Supplied local UnsolvedMath registry record: permanent ID \(701593\), exact source problem code \texttt{catalogue entry}, title \emph{Kourovka Notebook Problem 19.45}. The registry labels the problem \emph{open}; this is a supplied status, not an independently verified mathematical or literature assessment. The record was created and updated on 2026-10-16.
\item Attributed source: \emph{The Kourovka Notebook (Unsolved Problems in Group Theory)}, Problem 19.45, proposed by S. Kohl in 2018; the record identifies the 19th issue (2018) and the 21st issue edition (2026), \texttt{arXiv:1401.0300v48}. Its collection description names E. I. Khukhro and V. D. Mazurov as editors.
\item Kourovka Notebook Problem 17.57 is cited in the supplied statement for the definition of class transpositions. Its text is not supplied, and the reference has not been independently consulted.
\item The asserted sharpness example and finiteness claim are taken from the supplied statement of Problem 19.45. No separate supporting reference is supplied.
\end{itemize}
\end{document}
